\scanpage{388}
\mspage{112}
% archive sheet number: 387
$f_!$, en particulier des objets d'homologie $H_{\unsure{\bullet|0}}^{\mathrm{HOT}_W}(S,\unsure{\xi|\zeta})$ (ou simplement $H_{\bullet}^W(S,\xi)$) pour $\xi\in\mathrm{HOT}_W(e)=\mathrm{Hot}_W$, permet d'interpréter \unsure{homologiquement|homologiq\textsuperscript{mt}} la condition de $\mathrm{HOT}_W$-équivalence, comme signifiant que pour \unsure{tt} $\xi=\mathrm{hot}_W(\deleted{\unsure{I}}X)$ dans $\mathrm{Hot}_W$, la \deleted{\unsure{appli}} flèche canonique \unsure{de|dans} $\mathrm{Hot}_W$
\begingroup
\makeatletter
\tagsleft@true
\makeatother
\[
\tag*{$(\alpha)$}
H_{\bullet}^{\deleted{\unsure{D}}W}(S,\xi)\,\unsure{\relbar|\longrightarrow}\,H_{\bullet}^{W}(T,\xi)
\]
\endgroup
est iso. Or le calcul des $f_!$ (dans le cas de $p:S\relbar e$ et de $T\relbar e$) nous montre que les deux mbs de $(\alpha)$ ne sont autres \deleted{que $\unsure{S\times{}}$} (à isos près) que $\mathrm{hot}_W(X\times S)$ et $\mathrm{hot}_W(X\times T)$, la flèche étant induite par
\begingroup
\makeatletter
\tagsleft@true
\makeatother
\[
\mathrm{id}_X\times f:X\times S\longrightarrow X\times T .
\tag*{$(\alpha)$}
\]
\endgroup
Donc la condition envisagée signifie que cette flèche est un iso dans $\mathrm{Hot}_W$ pour tt $X$ dans $\mathrm{Cat}$. Prenant $X=e$, on voit que cela implique que $\mathrm{hot}_W(f)$ est iso, et cela suffit, vu que \unsure{la|le}
